\lecture{11}{KKT-условия и градиентный спуск}

\sourcevideo
    {https://www.youtube.com/playlist?list=PLOzRYVm0a65fhKDS8cYIC7YCuVGJ4FOJx}
    {Week 3: Lecture 11: KKT Conditions}

Условия Каруша--Куна--Таккера объединяют допустимость, стационарность
и согласование прямых и двойственных переменных. Для выпуклых задач
они дают проверяемый сертификат оптимальности. Во второй части лекции
начинается анализ градиентного спуска для выпуклых \(L\)-гладких
функций и выводится стандартный шаг \(1/L\).

\section{KKT-условия для задач с ограничениями}

Рассмотрим задачу
\[
    \begin{aligned}
        \minimize_{x\in\RR^n}
        \quad&
        f(x),
        \\
        \subjectto
        \quad&
        h_i(x)\le0,
        \qquad i=1,\ldots,m,
        \\
        &
        \ell_j(x)=0,
        \qquad j=1,\ldots,r.
    \end{aligned}
    \tag{\ensuremath{P}}
\]

Предполагается, что функции \(f\), \(h_i\) и \(\ell_j\)
дифференцируемы. Функция Лагранжа равна
\[
    \mathcal L(x,\lambda,\nu)
    =
    f(x)
    +
    \sum_{i=1}^{m}\lambda_i h_i(x)
    +
    \sum_{j=1}^{r}\nu_j\ell_j(x),
\]
где
\[
    \lambda\in\RR^m,
    \qquad
    \nu\in\RR^r.
\]

Переменные \(\lambda_i\) соответствуют ограничениям-неравенствам,
а \(\nu_j\) --- ограничениям-равенствам.



\begin{definition}[Условия Каруша--Куна--Таккера]
Тройка
\[
    (x^\star,\lambda^\star,\nu^\star)
\]
удовлетворяет условиям Каруша--Куна--Таккера, если выполнены четыре
группы условий.

Стационарность:
\[
    \nabla_x
    \mathcal L
    (x^\star,\lambda^\star,\nu^\star)
    =
    0.
\]

Прямая допустимость:
\[
    h_i(x^\star)\le0,
    \qquad
    \ell_j(x^\star)=0.
\]

Двойственная допустимость:
\[
    \lambda_i^\star\ge0.
\]

Дополняющая нежёсткость:
\[
    \lambda_i^\star h_i(x^\star)=0,
    \qquad i=1,\ldots,m.
\]
\end{definition}

Условие дополняющей нежёсткости утверждает, что для каждого
ограничения выполняется хотя бы одно из равенств
\[
    \lambda_i^\star=0
    \qquad\text{или}\qquad
    h_i(x^\star)=0.
\]

\begin{definition}[Активное ограничение]
Ограничение \(h_i(x)\le0\) называется активным в точке \(x^\star\),
если
\[
    h_i(x^\star)=0.
\]
Если
\[
    h_i(x^\star)<0,
\]
оно называется неактивным.
\end{definition}

Для неактивного ограничения дополняющая нежёсткость немедленно даёт
\[
    \lambda_i^\star=0.
\]
Следовательно, в условии стационарности участвуют только активные
ограничения с потенциально ненулевыми множителями.



В выпуклой постановке функции \(f\) и \(h_i\) выпуклы и
дифференцируемы, а ограничения-равенства аффинны:
\[
    \ell_j(x)=a_j^{\mathsf T}x-b_j.
\]

\begin{theorem}[KKT как критерий оптимальности]
Пусть задача \((P)\) выпукла и все входящие в неё функции
дифференцируемы. Любая тройка
\[
    (x^\star,\lambda^\star,\nu^\star),
\]
удовлетворяющая условиям KKT, задаёт глобальное решение прямой задачи
и оптимальные двойственные множители.

Обратно, пусть \(x^\star\) является решением прямой задачи и
выполнено условие Слейтера. Тогда существуют множители
\((\lambda^\star,\nu^\star)\), с которыми тройка
\((x^\star,\lambda^\star,\nu^\star)\) удовлетворяет условиям KKT.
\end{theorem}

\begin{proof}[Достаточность]
Поскольку \(\lambda^\star\ge0\), функция
\[
    x\longmapsto
    \mathcal L(x,\lambda^\star,\nu^\star)
\]
выпукла. Стационарность поэтому означает, что \(x^\star\) является
её глобальным минимумом. Следовательно,
\[
    g(\lambda^\star,\nu^\star)
    =
    \mathcal L
    (x^\star,\lambda^\star,\nu^\star).
\]

Из прямой допустимости и дополняющей нежёсткости
\[
\begin{aligned}
    \mathcal L
    (x^\star,\lambda^\star,\nu^\star)
    &=
    f(x^\star)
    +
    \sum_i
    \lambda_i^\star h_i(x^\star)
    +
    \sum_j
    \nu_j^\star\ell_j(x^\star)
    \\
    &=
    f(x^\star).
\end{aligned}
\]
Следовательно,
\[
    g(\lambda^\star,\nu^\star)
    =
    f(x^\star).
\]
Теперь используем слабую двойственность и допустимость
\(x^\star\):
\[
    g(\lambda^\star,\nu^\star)
    \le
    d^\star
    \le
    p^\star
    \le
    f(x^\star).
\]
Крайние величины равны, поэтому равенства выполняются во всей цепочке:
\[
    g(\lambda^\star,\nu^\star)
    =d^\star=p^\star=f(x^\star).
\]
\end{proof}

Точные KKT-условия дают сертификат оптимальности. Малые KKT-невязки
часто используются как критерий остановки, но сами по себе не
гарантируют близости к решению: для количественной оценки нужны
дополнительные условия регулярности или оценки ошибки.

На практике проверяются невязки:
\[
    \norm{
        \nabla_x
        \mathcal L(x,\lambda,\nu)
    },
\]
\[
    \positivepart{h_i(x)},
    \qquad
    \abs{\ell_j(x)},
    \qquad
    \positivepart{-\lambda_i},
\]
\[
    \abs{\lambda_i h_i(x)}.
\]

\section{Примеры KKT: квадратичная задача и SVM}

Рассмотрим
\[
    \begin{aligned}
        \minimize_{x\in\RR^n}
        \quad&
        \frac12x^{\mathsf T}Qx+c^{\mathsf T}x,
        \\
        \subjectto
        \quad&
        Ax=b,
    \end{aligned}
    \tag{\ensuremath{QP}}
\]
где
\[
    Q=Q^{\mathsf T}\succ0.
\]

Функция Лагранжа равна
\[
    \mathcal L(x,\nu)
    =
    \frac12x^{\mathsf T}Qx
    +
    c^{\mathsf T}x
    +
    \nu^{\mathsf T}(Ax-b).
\]

Ограничений-неравенств нет, поэтому условия двойственной допустимости
и дополняющей нежёсткости отсутствуют.

Стационарность по \(x\) даёт
\[
    Qx^\star+c+A^{\mathsf T}\nu^\star=0.
\]
Прямая допустимость имеет вид
\[
    Ax^\star=b.
\]

Эти равенства объединяются в систему
\[
    \begin{pmatrix}
        Q & A^{\mathsf T}
        \\
        A & 0
    \end{pmatrix}
    \begin{pmatrix}
        x^\star
        \\
        \nu^\star
    \end{pmatrix}
    =
    \begin{pmatrix}
        -c
        \\
        b
    \end{pmatrix}.
\]

Матрица
\[
    \begin{pmatrix}
        Q & A^{\mathsf T}
        \\
        A & 0
    \end{pmatrix}
\]
называется KKT-матрицей. При \(Q\succ0\) и полном строчном ранге
\(A\) она невырождена, поэтому решение \((x^\star,\nu^\star)\)
единственно.

\begin{remark}
Для простой квадратичной задачи KKT-система непосредственно даёт
решение. В более сложных задачах условия KKT обычно используются как
критерий остановки итерационного алгоритма.
\end{remark}



Для строго линейно разделимой выборки рассмотрим прямую задачу
\[
    \begin{aligned}
        \minimize_{w\in\RR^n,\ b\in\RR}
        \quad&
        \frac12\norm{w}^2,
        \\
        \subjectto
        \quad&
        y_i(w^{\mathsf T}x_i+b)\ge1,
        \qquad i=1,\ldots,m.
    \end{aligned}
\]

Запишем ограничения в стандартной форме:
\[
    h_i(w,b)
    =
    1-y_i(w^{\mathsf T}x_i+b)
    \le0.
\]

Функция Лагранжа равна
\[
    \mathcal L(w,b,\lambda)
    =
    \frac12\norm{w}^2
    +
    \sum_{i=1}^{m}
    \lambda_i
    \left[
        1-y_i(w^{\mathsf T}x_i+b)
    \right].
\]

Условия стационарности:
\[
    \nabla_w\mathcal L=0
    \quad\Longleftrightarrow\quad
    w^\star
    =
    \sum_{i=1}^{m}
    \lambda_i^\star y_i x_i,
\]
\[
    \frac{\partial\mathcal L}{\partial b}=0
    \quad\Longleftrightarrow\quad
    \sum_{i=1}^{m}\lambda_i^\star y_i=0.
\]

Прямая допустимость:
\[
    y_i(w^{\star\mathsf T}x_i+b^\star)\ge1.
\]

Двойственная допустимость:
\[
    \lambda_i^\star\ge0.
\]

Дополняющая нежёсткость:
\[
    \lambda_i^\star
    \left[
        1-y_i(w^{\star\mathsf T}x_i+b^\star)
    \right]
    =
    0.
\]

Если
\[
    \lambda_i^\star>0,
\]
то
\[
    y_i(w^{\star\mathsf T}x_i+b^\star)=1.
\]
Такие точки лежат на границе максимального отступа.

\begin{definition}[Опорный вектор]
Точку \(x_i\) будем называть опорным вектором двойственного
представления, если
\[
    \lambda_i^\star>0.
\]
\end{definition}

В вырожденном случае точка может лежать на границе отступа и при этом
иметь нулевой множитель. Поэтому геометрическое множество граничных
точек и носитель двойственного решения не обязаны совпадать.

Для любого опорного вектора
\[
    b^\star
    =
    y_i-w^{\star\mathsf T}x_i.
\]
На практике значение \(b^\star\) можно усреднить по нескольким
опорным векторам.

Точки с \(\lambda_i^\star=0\) не входят в представление
\[
    w^\star
    =
    \sum_i\lambda_i^\star y_i x_i.
\]
Если большинство множителей нулевые, решение имеет разреженное
представление через опорные объекты. Однако малый размер этого
множества не следует из KKT-условий автоматически.

\section{Геометрия градиентного спуска и импульс}

Рассмотрим задачу без ограничений
\[
    \minimize_{x\in\RR^n}f(x).
\]
Градиентный спуск имеет вид
\[
    x^{k+1}
    =
    x^k-\eta\nabla f(x^k).
\]

Градиент ортогонален линии уровня функции и указывает направление
наибольшего локального возрастания. Поэтому вектор
\[
    -\nabla f(x^k)
\]
задаёт направление наискорейшего локального убывания в евклидовой
норме.

Для плохо обусловленных квадратичных функций линии уровня сильно
вытянуты. В этом случае градиентный спуск может многократно пересекать
узкую долину, образуя зигзагообразную траекторию.



Для уменьшения колебаний можно учитывать направление предыдущего
шага:
\[
    x^{k+1}
    =
    x^k
    -
    \eta\nabla f(x^k)
    +
    \beta(x^k-x^{k-1}),
\]
где
\[
    \beta\in[0,1)
\]
--- параметр импульса.

При согласованной инициализации
\[
    x^k-x^{k-1}=-\eta v^k
\]
эквивалентная запись использует дополнительную переменную:
\[
    v^{k+1}
    =
    \beta v^k+\nabla f(x^k),
\]
\[
    x^{k+1}
    =
    x^k-\eta v^{k+1}.
\]

Предыдущая скорость сглаживает резкие изменения направления
градиента и может ускорить движение вдоль вытянутой долины.

\begin{remark}[Классификация метода]
В транскрипте метод с импульсом сопоставляется с методами второго
порядка. В строгой терминологии это метод первого порядка с памятью:
он использует градиенты, но не использует Гессиан. Дополнительная
стоимость состоит в хранении ещё одного вектора размерности \(n\).
\end{remark}

Для модели с миллиардами параметров дополнительный вектор состояния
может требовать значительной памяти. Поэтому скорость по числу
итераций необходимо сравнивать с фактической стоимостью одной
итерации и объёмом памяти.

\section{Лемма об убывании и шаг градиентного спуска}

Далее предполагается, что \(f\) выпукла и \(L\)-гладка:
\[
    \norm{\nabla f(x)-\nabla f(y)}
    \le
    L\norm{x-y}.
\]

\begin{lemma}[Лемма об убывании]
Для любой \(L\)-гладкой функции
\[
    f(y)
    \le
    f(x)
    +
    \inner{\nabla f(x)}{y-x}
    +
    \frac{L}{2}\norm{y-x}^2.
\]
\end{lemma}

\begin{proof}
Положим
\[
    \gamma(t)=x+t(y-x),
    \qquad t\in[0,1].
\]
Тогда
\[
\begin{aligned}
    f(y)-f(x)
    &=
    \int_0^1
    \inner{
        \nabla f(\gamma(t))
    }{
        y-x
    }
    \dd t
    \\
    &=
    \inner{\nabla f(x)}{y-x}
    \\
    &\quad+
    \int_0^1
    \inner{
        \nabla f(\gamma(t))-\nabla f(x)
    }{
        y-x
    }
    \dd t.
\end{aligned}
\]
По неравенству Коши--Буняковского и \(L\)-гладкости
\[
\begin{aligned}
    \inner{
        \nabla f(\gamma(t))-\nabla f(x)
    }{
        y-x
    }
    &\le
    Lt\norm{y-x}^2.
\end{aligned}
\]
Интегрирование даёт
\[
    f(y)-f(x)
    \le
    \inner{\nabla f(x)}{y-x}
    +
    \frac{L}{2}\norm{y-x}^2.
\]
\end{proof}

Правая часть является квадратичной верхней моделью функции около
точки \(x\).



Зафиксируем текущую точку \(x^k\). Из леммы об убывании
\[
\begin{aligned}
    f(y)
    \le{}&
    f(x^k)
    +
    \inner{\nabla f(x^k)}{y-x^k}
    +
    \frac{L}{2}\norm{y-x^k}^2.
\end{aligned}
\]

Обозначим
\[
    g^k=\nabla f(x^k).
\]
Дополняя квадрат, получаем
\[
\begin{aligned}
    f(y)
    \le{}&
    f(x^k)
    -
    \frac{1}{2L}\norm{g^k}^2
    \\
    &+
    \frac{L}{2}
    \norm{
        y-
        \left(
            x^k-\frac1L g^k
        \right)
    }^2.
\end{aligned}
\]

Правая часть минимальна при
\[
    y
    =
    x^k-\frac1L\nabla f(x^k).
\]
Поэтому естественный шаг градиентного спуска равен
\[
    \eta=\frac1L,
\]
а обновление имеет вид
\[
    x^{k+1}
    =
    x^k-\frac1L\nabla f(x^k).
\]

\begin{theorem}[Гарантированное убывание]
Для шага \(\eta=1/L\)
\[
    f(x^{k+1})
    \le
    f(x^k)
    -
    \frac{1}{2L}
    \norm{\nabla f(x^k)}^2.
\]
\end{theorem}

Если
\[
    \nabla f(x^k)\ne0,
\]
то
\[
    f(x^{k+1})<f(x^k).
\]
Следовательно, последовательность значений целевой функции
невозрастает.

\begin{remark}[Исправление коэффициента]
При дополнении квадрата коэффициент перед нормой градиента равен
\[
    \frac{1}{2L},
\]
а не \(1/L\). Эта арифметическая неточность появляется в устном
выводе и затем исправляется.
\end{remark}



\begin{algorithm}[H]
\caption{Градиентный спуск для \(L\)-гладкой функции}
\label{alg:lecture11-gradient-descent}
\begin{algorithmic}[1]

\Require
функция \(f\), константа гладкости \(L>0\),
начальная точка \(x^0\)

\For{\(k=0,1,2,\ldots\)}

    \State
    \(g^k\gets\nabla f(x^k)\)

    \If{\(\norm{g^k}\) достаточно мала}

        \State
        \Return \(x^k\)

    \EndIf

    \State
    \(x^{k+1}\gets x^k-\dfrac1L g^k\)

\EndFor

\end{algorithmic}
\end{algorithm}

Если \(L\) неизвестна, фиксированный шаг \(1/L\) непосредственно
использовать нельзя. На практике шаг выбирают экспериментально или
определяют посредством поиска по направлению. Эти процедуры
рассматриваются отдельно.

\section{Скорость для выпуклой гладкой функции}

Пусть \(x^\star\) --- глобальный минимум выпуклой \(L\)-гладкой
функции.

\begin{theorem}[Сублинейная сходимость]
Для градиентного спуска
\[
    x^{k+1}
    =
    x^k-\frac1L\nabla f(x^k)
\]
выполняется
\[
    f(x^k)-f(x^\star)
    \le
    \frac{
        L\norm{x^0-x^\star}^2
    }{
        2k
    },
    \qquad k\ge1.
\]
\end{theorem}

\begin{proof}
Обозначим \(g^k=\nabla f(x^k)\). По выпуклости
\[
    f(x^k)-f(x^\star)
    \le
    \inner{g^k}{x^k-x^\star}.
\]
По гарантированному убыванию
\[
    f(x^{k+1})
    \le
    f(x^k)-\frac{1}{2L}\norm{g^k}^2.
\]
Следовательно,
\[
\begin{aligned}
    f(x^{k+1})-f(x^\star)
    \le{}&
    \inner{g^k}{x^k-x^\star}
    -
    \frac{1}{2L}\norm{g^k}^2
    \\
    ={}&
    \frac{L}{2}
    \left(
        \norm{x^k-x^\star}^2
        -
        \norm{x^{k+1}-x^\star}^2
    \right).
\end{aligned}
\]
Суммируя по \(k=0,\ldots,N-1\), получаем
\[
    \sum_{k=0}^{N-1}
    \bigl(
        f(x^{k+1})-f(x^\star)
    \bigr)
    \le
    \frac{L}{2}
    \norm{x^0-x^\star}^2.
\]
Поскольку \(f(x^k)\) невозрастает,
\[
    N\bigl(f(x^N)-f(x^\star)\bigr)
    \le
    \frac{L}{2}
    \norm{x^0-x^\star}^2.
\]
\end{proof}

Итак,
\[
    f(x^k)-f^\star
    =
    \BigO\left(\frac1k\right).
\]

Чтобы гарантировать
\[
    f(x^k)-f^\star\le\varepsilon,
\]
достаточно
\[
    k
    \ge
    \frac{
        L\norm{x^0-x^\star}^2
    }{
        2\varepsilon
    }.
\]
Итерационная сложность составляет
\[
    \BigO\left(\frac1\varepsilon\right).
\]

\section{Асимптотические обозначения и порядок сходимости}

\begin{definition}[Большое \(O\)]
Запись
\[
    a_k=\BigO(b_k)
\]
означает, что существуют \(C>0\) и \(k_0\), такие что
\[
    \abs{a_k}
    \le
    C\abs{b_k}
\]
для всех \(k\ge k_0\).
\end{definition}

Например,
\[
    k^2-2k+3k^3
    =
    \BigO(k^3).
\]

Оценка
\[
    f(x^k)-f^\star
    =
    \BigO(k^{-1})
\]
описывает скорость уменьшения ошибки по значению функции, а не
обязательно скорость уменьшения расстояния
\(\norm{x^k-x^\star}\).



Пусть
\[
    x^k\to x^\star,
    \qquad
    e_k=\norm{x^k-x^\star},
\]
и \(e_k>0\) для всех достаточно больших \(k\).

\begin{definition}[Порядок сходимости]
Последовательность имеет порядок сходимости \(q\ge1\), если
существует \(\rho\in(0,\infty)\), для которого
\[
    \lim_{k\to\infty}
    \frac{e_{k+1}}{e_k^q}
    =
    \rho.
\]
Число \(\rho\) называется асимптотической константой ошибки.
\end{definition}

При \(q=1\) говорят о линейной сходимости, если \(\rho<1\).
При \(q=2\) сходимость называется квадратичной.

Оценку
\[
    f(x^k)-f^\star=\BigO(k^{-1})
\]
называют сублинейной, поскольку она не задаёт геометрического
убывания ошибки. Из одной оценки большого \(O\) нельзя заключить, что
отношение последовательных ошибок обязательно стремится к единице.

Сильная выпуклость позволит получить более быструю линейную
сходимость градиентного метода. При дополнительных предположениях и подходящем выборе параметров
методы с импульсом позволяют улучшить зависимость скорости от числа
обусловленности.

\begin{lecturesummary}
Условия Каруша--Куна--Таккера состоят из стационарности, прямой и
двойственной допустимости и дополняющей нежёсткости. Для выпуклой
задачи они достаточны для глобальной оптимальности, а при условии
регулярности также необходимы. В квадратичной задаче с равенствами
KKT-условия образуют линейную систему с седловой KKT-матрицей. Для
SVM положительные двойственные множители выделяют опорные векторы и
позволяют восстановить смещение разделяющей гиперплоскости.

Для \(L\)-гладкой функции выполняется лемма об убывании
\[
    f(y)
    \le
    f(x)
    +
    \inner{\nabla f(x)}{y-x}
    +
    \frac{L}{2}\norm{y-x}^2.
\]
Минимизация этой верхней модели приводит к шагу
\[
    x^{k+1}
    =
    x^k-\frac1L\nabla f(x^k)
\]
и оценке
\[
    f(x^{k+1})
    \le
    f(x^k)
    -
    \frac{1}{2L}
    \norm{\nabla f(x^k)}^2.
\]
Для выпуклой \(L\)-гладкой функции градиентный спуск имеет скорость
\[
    f(x^k)-f^\star
    =
    \BigO\left(\frac1k\right)
\]
и требует порядка \(1/\varepsilon\) итераций для достижения ошибки
\(\varepsilon\) по значению функции.
\end{lecturesummary}